%%%%%%%%%%%%%%%%%%%%%%%%%%%%%%%%%%%%%%%%
%        Author: Bernard Lampe
%   From Primitive to Exploit lecture material
%%%%%%%%%%%%%%%%%%%%%%%%%%%%%%%%%%%%%%%%

% import template
\documentclass[12pt]{Training}

% import packages
\usepackage[utf8]{inputenc}
\usepackage[T1]{fontenc}
\usepackage{mathpazo}
\usepackage{graphicx}
\usepackage{booktabs}
\usepackage{listings}
\usepackage{enumerate}
\usepackage{xcolor}

%----------------------------------------------------------------------------------------

% set document info
\title{From Primitive to Exploit Lecture Material}
\author{Dr. Bernard Lampe}
\class{Introduction to Vulnerability Research}
\date{November 21st, 2019}

%----------------------------------------------------------------------------------------

% set code listing style
\lstset{
    basicstyle=\ttfamily\small,
    numbers=left,
    tabsize=2,
    keywordstyle=\color{blue}\ttfamily,
    stringstyle=\color{red}\ttfamily,
    commentstyle=\color{green}\ttfamily,
    morecomment=[l][\color{magenta}]{\#}
}

%----------------------------------------------------------------------------------------

\begin{document}

% print title
\maketitle

%----------------------------------------------------------------------------------------

\section*{Introduction}
\begin{enumerate}
    \item A memory-corruption finding needs to be paired with a goal and technique. This is where vulnerability research transfers to exploit engineering.
    \item Scope statement: this is a \emph{recognition-level} treatment --- how the technique families work, what they need, why each stops working. Building a real chain end-to-end is an exploitation course in itself.
\end{enumerate}

\section*{Control-Flow Reuse: ROP and JOP}
\begin{enumerate}
    \item The NX/DEP mitigation and its bypass family. Premise: with the stack non-executable, reuse existing executable bytes instead of injecting code.
    \item \textbf{Ret2libc}: the trivial chain --- overwrite the return address to jump to \lstinline{system("sh")} (or the platform equivalent) with argument registers/stack pre-filled per the calling convention. One function, no gadget search: the first technique to try against any no-canary stack overflow.
    \item \textbf{ROP} (Return-Oriented Programming): chain short ``gadget'' snippets ending in \lstinline{ret}; each gadget performs a compute step (register move, add, dereference-load). The chain lives on the stack; each return consumes the next gadget address. Bootstrap requirement: a controlled return address and controlled stack contents behind it.
    \begin{enumerate}
        \item Gadget discovery: scan the binary (and loaded libraries) for \lstinline{ret}-terminated fragments --- ROPgadget/ropper class tools; after PIE defeat you can mix libraries.
        \item Calling a syscall on x86\_64 needs \lstinline{rax}=NR, \lstinline{rdi/rsi/rdx} args: typical chains are \emph{load gadget} $\to$ \emph{pop gadgets} $\to$ \lstinline{syscall}; the gadget census of the target decides feasibility.
        \item Chain fragility: stack alignment (x86 \lstinline{movaps} traps on unaligned), NULL bytes limiting address choice, and the count of useful gadgets --- small static targets may have too few, requiring ret2libc or more bugs to get there.
    \end{enumerate}
    \item \textbf{JOP} (Jump-Oriented): same idea with \lstinline{jmp} dispatch through a register with a dispatcher gadget; survives CFI-in-a-can (ret-only policies) and is the standard answer where the toolchain removed \lstinline{ret}-gadgets specifically.
    \item \textbf{Stack pivot}: when the initial control is a single indirect call/branch (function-pointer overwrite) rather than a return, first redirect rsp into attacker data (\lstinline{xchg esp, eax} class gadgets), \emph{then} run a ROP chain. Prerequisite chain: controlled call $\to$ pivot $\to$ chain --- why function-pointer hijacks pair with pivots in every writeup.
    \item The variants that recur, and the condition each needs:
    \begin{enumerate}
        \item \textbf{ret2libc} with argument control: the same as the trivial chain, but the argument registers (\lstinline{rdi} on x86-64, \lstinline{x0} on ARM64) must be populated first --- so it needs a \lstinline{pop rdi; ret} gadget and a pointer to \lstinline{"/bin/sh"}. Without argument control it is \lstinline{system} called with whatever happens to be in the register.
        \item \textbf{ret2plt}/\textbf{ret2dlresolve}: call an existing PLT stub directly (no leak needed for the call target), or make the dynamic resolver resolve an attacker-named symbol. ret2dlresolve needs control of the relocation structures on the stack and is the no-libc-leak answer on a PIE binary.
        \item \textbf{ret2csu}: the \lstinline{__libc_csu_init} gadget populates several argument registers at once from the stack, which is why it is the standard first gadget when no \lstinline{pop} gadgets exist.
        \item \textbf{SROP} (sigreturn-oriented programming): the attacker forges a \lstinline{sigcontext} on the stack and triggers \lstinline{rt_sigreturn}, which loads \emph{all} registers at once. It needs only a \lstinline{syscall} gadget and a controlled stack, which is why it is the standard answer when the gadget census is too small for a normal chain.
        \item \textbf{BROP}/\textbf{JIT-ROP}: when the binary is not available (remote, or a browser target), leak the code at runtime and build the chain on the fly. This is dynamic tooling applied to exploit construction.
    \end{enumerate}
    \item The modern blocks and their bypasses: \textbf{CET} (shadow stack, \lstinline{endbr64}) blocks a plain \lstinline{ret}-redirect unless a legitimate \lstinline{endbr64} landing pad exists, which pushes toward data-only attacks or \lstinline{call}-oriented chains; \textbf{ARM64 PAC} (\lstinline{paciasp}/\lstinline{autiasp}) blocks overwriting a saved return address unless the PAC value can be forged or the signed pointer replayed; \textbf{MTE} blocks tag-mismatched accesses. Read the target's \lstinline{/proc/cpuinfo} and compiler flags before choosing a technique.
    \item Tooling: \lstinline{ROPgadget}/\lstinline{ropper}/\lstinline{rp++} for the gadget census, \lstinline{pwntools ROP} for chain assembly and \lstinline{one_gadget} for the \lstinline{system}-and-nothing-else case. The chain is assembled, then tested, then re-assembled; the debugger is the loop.
\end{enumerate}

\section*{Infoleak: Defeating ASLR Before Any Chain}
\begin{enumerate}
    \item Every ROP-family chain needs absolute addresses. ASLR randomizes them per-boot; the chain therefore needs a leak first --- this is why ``bug 1 = leak, bug 2 = control'' is the standard exploitation shape.
    \item Leak primitive families:
    \begin{enumerate}
        \item Reflected OOB reads whose output length/reflection is attacker-influenced (format string reads, HeartBleed-class over-reads, unbounded \lstinline{snprintf} echoes).
        \item Uninitialized-memory reads returning stale heap/stack content (the uninitialized-use bug class's flip side: attacker sprays pointers, then triggers a use that echoes them).
        \item Side channels: timing (cache sets, branch predictors), error messages, crash-adjacent dumps, and pointer-value echo through logs.
    \end{enumerate}
    \item What to leak: one pointer per randomized region (libc for gadgets, stack for pivots, heap for groom-relative addresses). Then the ASLR row of the mitigation table collapses for that session.
    \item Assessing a leak: bytes/attempt, retry budget (does the connection survive the read?), and whether the leak's target address survives between sessions (per-boot vs per-process randomization).
    \item The leak sources that recur, by primitive:
    \begin{enumerate}
        \item A reflected over-read: \lstinline{printf} with an attacker-controlled \lstinline{%p}/\lstinline{%s} argument, an unbounded copy-out, a heartbeat-style over-read (CVE-2014-0160). The bug is a read whose length or target the attacker controls.
        \item Uninitialized memory returned to the attacker: a partially filled struct copied out, a freed-then-reused buffer echoed back (the uninitialized-use bug class).
        \item Allocator metadata: a \lstinline{main_arena} pointer or an \lstinline{unsorted bin} \lstinline{fd}/\lstinline{bk} read from a corrupted chunk, or a \lstinline{tcache} entry pointing into \lstinline{libc}. This is the standard glibc base leak and needs only a relative read.
        \item A stack or \lstinline{environ} pointer: gives the stack base, which is what a \lstinline{SROP} or a stack pivot chain needs.
        \item The binary's own base: a partial pointer or a mapping address (\lstinline{/proc/self/maps} if readable) defeats PIE; on a fully PIE target the leaked \lstinline{libc} address plus the known link-time offset is often enough.
    \end{enumerate}
    \item The leak consumers, one per randomized region: a \lstinline{libc} leak for the gadgets and \lstinline{system}, a stack leak for a pivot or \lstinline{SROP}, a heap leak for the groom's relative addresses, and a binary leak for PIE. One leak usually yields the rest, because the libraries' relative offsets are fixed at link time.
    \item Tooling: \lstinline{pwntools} \lstinline{LibcSearcher} and the \lstinline{libc-database} identify a libc from one leaked symbol's offset; \lstinline{one_gadget} maps a libc build to a single-call chain; \lstinline{ASLR} on 32-bit is brute-forceable in a loop, which is why the retry budget above matters.
\end{enumerate}

\section*{Allocator Grooming: Making the Primitive Land}
\begin{enumerate}
    \item Heap primitives (sizes/lengths/double-free/UAF) convert into control only if the corrupted allocator state serves an attacker-shaped allocation next. Grooming is the heap-layout side of every UAF/overflow exploit: the question any heap primitive raises is which allocation the corrupted state serves next.
    \item glibc's allocator internals (the model built by walking a small program's allocations; versioned --- verify against the target's glibc):
    \begin{enumerate}
        \item Small allocations go to \textbf{tcache} (per-thread, size-binned singly-linked lists; first-fit; a double-free of one tcache entry yields two allocations of the same address --- 2.29+ prints ``detected double free'' on an exact same-size repeat, so shapes must oblige).
        \item \textbf{fastbins}, \textbf{unsorted bin}, \textbf{small/large bins}: the free-list machinery behind bigger chunks; metadata (fd/bk pointers, size fields) is what a controlled overflow re-writes to force ``next allocation serves my target''.
        \item The classic shapes: overwrite the size field to overlap chunks (House of Elephant/Force family), corrupt tcache fd to allocate over a function pointer, and unsorted-bin attacks to leak a libc address directly.
        \item Conservatism rules: every allocator hardening update (safe-linking in 2.32, tcache key double-free guard in 2.29/2.34, fd-pointer mangling) changed the toolbox --- check the \emph{target's} libc version first.
    \end{enumerate}
    \item Kernel slab (SLUB) parallels: per-size freelist in page-sized slabs; same grooming logic with different primitives (spray primitives such as msg\_msg, pipe buffers, add\_key); the cross-cache variant reclaims a freed slab page for a different cache --- the kernel form of ``groom until the type confusion is served''.
    \item Practical grooming loop: break on \lstinline{malloc/free}, record a transcript of sizes/order under the target's own command sequence, replay onto your exploit's predicted sequence, confirm the landing allocation receives your sprayed data (removing the guesswork one step at a time).
    \item The glibc facts that decide every groom, in the order an exploit needs them:
    \begin{enumerate}
        \item \textbf{tcache} (2.26+): per-thread, one bin per size, last-in-first-out, seven entries per bin by default. A free goes to tcache before any other bin, which is why a tcache-sized primitive is the most reliable one. The tunables (\lstinline{glibc.malloc.tcache_count}, \lstinline{glibc.malloc.tcache_max}) change the arithmetic, so read them.
        \item \textbf{fastbins}: small chunks with a LIFO singly-linked list; a double free is caught only for the exact same size (\lstinline{fasttop} check), so a double free of two sizes is the classic bypass.
        \item \textbf{unsorted/small/large bins}: freed chunks that are coalesced and sorted; their \lstinline{fd}/\lstinline{bk} pointers are the leak source and the \lstinline{unlink} target.
        \item \textbf{safe-linking} (2.32+): the freelist \lstinline{fd} is stored as \lstinline{(ptr >> 12) ^ pos}, so an attacker who overwrites it must know the chunk's own address --- a heap leak is now a precondition for the classic tcache-poisoning groom.
        \item \textbf{the tcache key} (2.29+): a freed tcache chunk stores a per-thread key, and \lstinline{free} refuses a second free of the same chunk (``\lstinline{free(): double free detected}''). The workaround is to overwrite the key, which itself needs a write primitive.
        \item \textbf{the mmap threshold}: allocations above it are served by \lstinline{mmap}, not the main arena, which changes which groom applies (\lstinline{M_MMAP_THRESHOLD}, \lstinline{glibc.malloc.mmap_threshold}).
    \end{enumerate}
    \item The named techniques, and the primitive each needs: \lstinline{house of force} (overwrite the top chunk size), \lstinline{house of spirit} (free a forged chunk pointer), \lstinline{house of einherjar} (forge a chunk's \lstinline{prev_size}), \lstinline{unsafe unlink} (a chunk whose \lstinline{fd}/\lstinline{bk} pass the check), \lstinline{tcache poisoning} (overwrite a tcache \lstinline{fd} to return an arbitrary address), and the \lstinline{_IO_FILE} family (\lstinline{FSOP}, \lstinline{house of apple}) which overwrites a file structure's vtable. Each is the allocator-metadata corruption pattern, assessed against the target's glibc version.
    \item The kernel parallel is the same logic in the SLUB allocator: a per-cache freelist in page-sized slabs, with \lstinline{CONFIG_SLAB_FREELIST_HARDENED} and \lstinline{_RANDOM} as the \lstinline{safe-linking} analogue, and cross-cache reclaim as the ``reuse a freed page for a different object type'' technique.
\end{enumerate}

\section*{Bypass Puzzle: Working a Real Chain}
\begin{enumerate}
    \item Worked example: the \lstinline{heap.c} string-store program included in this module (a UAF on double DELETE: the slot pointer stays live after \lstinline{free()}, so the next same-size CREATE lands on the freed chunk). With ASan the double-DELETE is caught immediately (quarantine); with plain glibc the UAF gives a freed-but-pointer-live allocation --- groom with same-size CREATEs so the re-CREATE lands on the stale slot, then FIND-after-free yields a type-confused read primitive. Each mitigation changes what the chain needs; walk each row below against this one bug.
    \item For each finding, write the requirements table: bug class $\times$ active mitigation $\times$ primitive ($=$ chain steps $\times$ required preconditions). The report's ``exploitability statement'' is exactly this table read aloud.
    \item The same table, filled in for the \lstinline{heap.c} double-DELETE (the UAF above), which is the worked example every report should imitate:
\end{enumerate}

\begin{center}
\footnotesize
\begin{tabular}{@{}p{3.2cm}|p{5.4cm}|p{5.6cm}@{}}
\toprule
\textbf{Mitigation} & \textbf{Effect on this bug} & \textbf{What the chain needs} \\
\midrule
ASan build & double-DELETE caught at once (quarantine) & nothing: the bug is not exploitable under ASan \\
glibc tcache (2.29+) & double free detected by the key check & overwrite the tcache key first, or free two different sizes \\
glibc safe-linking (2.32+) & tcache \lstinline|fd| is mangled by the chunk address & a heap leak before the poisoning groom \\
NX / DEP & no shellcode on the stack & a ROP or \lstinline|ret2libc| chain after the leak \\
ASLR / PIE & no fixed addresses & a \lstinline|libc| leak (unsorted-bin or \lstinline|_IO_FILE|) \\
FORTIFY\_SOURCE & checked \lstinline|strcmp|/\lstinline|memcpy| paths & use an unchecked path, or a controlled-length read \\
\bottomrule
\end{tabular}
\end{center}

\begin{enumerate}
    \item Reading the table: each row is a step in the chain, not a verdict. The exploit is the composition of the rows; the report is the table plus the one-sentence root cause.
    \item The general rule, which is the whole module: a primitive grants what the \emph{weakest} active mitigation lets it do. An arbitrary write on a fully-mitigated target needs a leak plus a groom plus a chain; the same write on an old target needs a single \lstinline{ret2libc}.
\end{enumerate}

\section*{Conclusion}
The bypass column is not magic: it is four techniques (code reuse, leaks, grooming, pivots) paired to bug classes by need. An auditor who cannot build the pairing systematically overstates what a finding grants; an exploit developer who skips the audit discipline rebuilds the same knowledge at higher cost.

The discipline to take away: for every finding, name the primitive, name the active mitigation, and name the next step. ``I have a UAF'' is not a finding; ``I have a same-size UAF on glibc 2.31, the tcache double-free check is the only obstacle, and I need a heap leak for safe-linking'' is. The three steps of every chain are the three sections of this module: get the addresses (infoleak), make the corruption land (grooming), then turn the landing into control (code reuse or a data-only write).

\section*{References}
\begin{enumerate}
	\item Return-oriented programming: original papers by Shacham (2007) and Checkoway et al (2010, JOP)
	\item pwntools: docs.pwntools.com (gadget search and chain assembly support)
	\item ROPgadget / ropper / rp++ (gadget census tools)
	\item glibc malloc internals, Marshall (Phrack ``Malloc Des Maleficarum'') and the glibc malloc wiki writeups
	\item how2heap: github.com/shellphish/how2heap (allocator technique drills)
	\item The Shellcoder's Handbook (2nd ed) --- mitigation-era exploitation chapters
	\item SROP: Bosman \& Bos, ``Framing Signals --- a Return to Portable Shellcode'' (2014)
	\item Data-only and counter-CFI: Hu et al, ``Data-Oriented Programming'' (2022) and Schuster et al, ``Counterfeit Object-Oriented Programming'' (2015)
	\item Blind and JIT-ROP: Bittau et al (2014) and Snow et al (2013)
	\item libc-database and one\_gadget (libc.rip, one-gadget)
\end{enumerate}

\end{document}
